\section{Diagnostic Protocol and Experimental Setup}

The diagnostic protocol treats an intervention-trigger claim as a cost-sensitive human-attention allocation claim rather than as a simple learned-trigger-versus-random comparison. A trigger is not credited for being selective merely because its nominal budget differs from the baseline budget; it must be compared with random intervention policies whose realized human-step costs bracket the learned trigger. A trigger is also not credited from a single favorable checkpoint; selected checkpoints must be evaluated with repeated episodes and reported with uncertainty. Finally, aggregate success and cost are not sufficient when a trigger fails, because they do not explain how attention was spent. The protocol therefore combines cost-matched random families, repeated checkpoint evaluation, and trace-level intervention diagnostics as three validity gates.

The first validity gate is cost matching. Cost matching is performed against realized best-step human cost, not only against the nominal intervention budget assigned at launch. For each learned trigger under test, the random family is swept across nearby budgets so that at least one random point falls below and one falls near or above the trigger's realized cost when feasible. This converts the comparison from "learned trigger versus a random baseline" into a local success-cost frontier. In the Lift case study, this step is decisive: the lower-budget random b350 condition dominates LV-VoI scale3, so the protocol rejects the trigger-superiority interpretation before any repair is expanded.

The second validity gate is repeated checkpoint evaluation. This gate reduces dependence on a single final policy or a single restored-best evaluation. For paper-facing Lift claims, selected checkpoints are evaluated with repeated rollouts and reported as success counts, success proportions, Wilson confidence intervals, and realized human-step cost. Wilson intervals are used because they behave more reliably than simple normal approximations for finite binomial samples, especially near high success rates \cite{wilson1927score}. The protocol also keeps best-step and final-policy accounting separate so that a policy's apparent success is not confused with late training collapse or checkpoint-selection noise.

The third validity gate is trace-level diagnosis. Trace diagnostics are used after a trigger fails a cost-matched comparison because a dominated aggregate result does not reveal whether the trigger spent attention too early, too often, in the wrong geometry, or under miscalibrated scores. Each intervention start records the training step, the fraction of budget already spent, gripper-to-cube geometry, end-effector and object height, trigger score, and predicted failure probability when those quantities are available. These traces make mechanism hypotheses testable. For example, R023 shows that LV-VoI does not simply fail by asking far from the cube; it starts closer to the cube than random on average while spending many more intervention starts. R024 then shows that suppressing some low-score mid-run starts is mechanically possible but insufficient to recover the success-cost frontier.

The current case study uses robosuite Lift and Stack as simulated robotic manipulation environments \cite{robosuite2020}. The learner follows an off-policy SAC-style setup \cite{haarnoja2018sac,sac2018applications} with demonstration and intervention data motivated by offline-to-online robot learning and demonstration-learning practice \cite{ball2023rlpd,mandlekar2021robomimic}. The evaluated intervention strategies include no intervention, a random-budget family, LV-VoI scale3, a minimum-disagreement LV-VoI repair, and a score-floor LV-VoI repair. Lift supplies the main cost-matched diagnostic evidence; Stack is retained as boundary evidence showing that the current intervention mechanism does not yet transfer as a positive method result.

The intervention source is a scripted privileged-state oracle, not a real teleoperator. This design makes intervention timing reproducible and cheap enough to diagnose across seeds, budgets, and checkpoints, but it limits claims about real human workload, interface latency, fatigue, and embodiment transfer. Accordingly, all numerical claims in this paper are routed through the evidence registry and registry-generated R036 tables rather than copied manually from raw logs. Table~\ref{tab:protocol_checklist} summarizes the reporting gates that follow from this diagnostic setup.

\begin{table*}[t]
\centering
\caption{Diagnostic protocol checklist for HIL-RL intervention-trigger claims. Each gate is motivated by a failure mode surfaced in the Lift/Stack case study.}
\label{tab:protocol_checklist}
\begin{tabular}{p{0.07\textwidth}p{0.19\textwidth}p{0.32\textwidth}p{0.33\textwidth}}
\toprule
Gate & Requirement & What to report & Why it matters \\
\midrule
G1 & Cost-matched random family & At least one random budget below and above the method's realized human cost; report realized best-step human cost, not only nominal budget. & Prevents a learned trigger from appearing efficient only because the random baseline was mismatched. \\
G2 & Repeated checkpoint evaluation & Evaluate selected checkpoints with repeated episodes; report binomial confidence intervals and seed variability. & Single restored or final evaluations can change rankings in noisy RL settings. \\
G3 & Best-step and final-policy accounting & Report raw final policy, restored best policy, and best-step human cost. & Separates late collapse/checkpoint selection from true intervention efficiency. \\
G4 & Trace-level intervention diagnostics & Log intervention starts with timing, budget fraction, task geometry, score, and p\_fail when available. & Explains how a trigger spends budget and tests mechanism hypotheses. \\
G5 & Same-seed stop gate & Before expanding a repaired trigger, compare it against same-seed cost-matched random using repeated evaluation. & Stops weak heuristic repairs before expensive expansion and prevents over-tuning. \\
\bottomrule
\end{tabular}
\end{table*}

